\documentclass[11pt]{amsart}
\usepackage[T1]{fontenc}
\usepackage{lmodern,microtype,amsmath,amssymb,amsthm,mathtools,enumitem,booktabs}
\usepackage[margin=1.1in]{geometry}
\usepackage[hidelinks]{hyperref}
\usepackage{xurl}

\newtheorem{theorem}{Theorem}[section]
\newtheorem{lemma}[theorem]{Lemma}
\newtheorem{corollary}[theorem]{Corollary}
\newtheorem{proposition}[theorem]{Proposition}
\newtheorem{conjecture}[theorem]{Conjecture}
\theoremstyle{remark}
\newtheorem{remark}[theorem]{Remark}
\numberwithin{equation}{section}

\newcommand{\Ent}{\operatorname{Ent}}
\newcommand{\tr}{\operatorname{tr}}
\newcommand{\adj}{\operatorname{adj}}

\hypersetup{pdftitle={Concavity of determinantal entropy on two points},pdfauthor={Yitzchak Shmalo}}

\title[Concavity of determinantal entropy on two points]{Concavity of the entropy of determinantal measures on two points, and the second variation in general}
\author{Yitzchak Shmalo}
\address{Einstein Institute of Mathematics, The Hebrew University of Jerusalem, Jerusalem, Israel}
\email{yitzchak.shmalo@gmail.com}
\date{6 November 2026}
\subjclass[2020]{60G55, 94A17, 15A15, 15A45}
\keywords{determinantal probability measure, entropy, concavity, principal minors}

\begin{document}
\begin{abstract}
For a positive contraction $Q$ on $\ell^2(E)$ with $E$ finite, let $P_Q$ be the determinantal probability measure on $2^E$ and $\Ent(Q)$ its Shannon entropy. Lyons conjectured in 2003 that $\Ent$ is concave in $Q$. We prove the conjecture when $|E|=2$. For general $E$ we write the second derivative of $\Ent$ along a line as a negative quadratic form in the derivatives of the principal minors of $Q$ plus a sum, over subsets $T$ with $|T|\ge2$, of the second derivatives of the minors $\det Q_T$ weighted by M\"obius sums of $\log P_Q$. For two-element $T$ these weights are nonnegative by the Hadamard--Fischer inequality, which is what makes the case $|E|=2$ work. We also report numerical evidence for $|E|=3,4$: the second derivative is negative away from scalar kernels $cI$, where it vanishes in the off-diagonal directions.
\end{abstract}
\maketitle

\section{Introduction}\label{sec:intro}

Let $E$ be a finite set and $Q$ a positive contraction on $\ell^2(E)$, that is, a Hermitian matrix with $0\le Q\le I$. The determinantal probability measure $P_Q$ on $2^E$ is determined by $P_Q[\,A\subseteq Y\,]=\det Q_A$ for $A\subseteq E$, where $Q_A$ is the principal submatrix indexed by $A$ \cite{Ly03}. Its entropy is
\[
\Ent(Q)\;=\;-\sum_{S\subseteq E}P_Q[Y=S]\log P_Q[Y=S].
\]
In Section~9 of \cite{Ly03}, Lyons asked whether $\Ent$ is concave on the convex set $\{0\le Q\le I\}$ and stated the following conjecture, supported by numerical calculation.

\begin{conjecture}[Lyons {\cite[Conjecture~9.2]{Ly03}}]\label{conj}
For positive contractions $Q_1,Q_2$ on $\ell^2(E)$, $\Ent\bigl((Q_1+Q_2)/2\bigr)\ge\bigl(\Ent(Q_1)+\Ent(Q_2)\bigr)/2$.
\end{conjecture}

The conjecture is restated as Conjecture~2.6 of Lyons's survey \cite{Ly14}. By finite restrictions it implies Conjecture~5.3 there: the entropy of the stationary determinantal process $P_f$ with symbol $f$, studied by Lyons and Steif \cite{LS03}, is a concave function of $f$. The map $Q\mapsto P_Q$ is polynomial of degree $|E|$, so concavity of $\Ent$ is not inherited from the concavity of $x\mapsto -x\log x$.

\begin{theorem}\label{thm:main}
Conjecture~\ref{conj} holds when $|E|=2$.
\end{theorem}

Section~\ref{sec:two} proves Theorem~\ref{thm:main}. Section~\ref{sec:general} gives the second variation of $\Ent$ for arbitrary finite $E$ (Proposition~\ref{prop:second}) and explains which part of the two-point argument survives. Section~\ref{sec:numerics} records numerical evidence for $|E|=3$ and $4$.

\section{Two points}\label{sec:two}

Let $E=\{1,2\}$ and write $Q=\begin{psmallmatrix}a&b\\\bar b&d\end{psmallmatrix}$. The four probabilities are
\begin{equation}\label{eq:p}
p_{12}=\det Q,\qquad p_1=a-\det Q,\qquad p_2=d-\det Q,\qquad p_0=1-a-d+\det Q,
\end{equation}
where $p_S=P_Q[Y=S]$ and $p_0=P_Q[Y=\emptyset]$. A direct computation gives
\begin{equation}\label{eq:cov}
p_1p_2-p_0p_{12}=|b|^2 .
\end{equation}
Since $\Ent$ is continuous on the compact set $\{0\le Q\le I\}$, it suffices to prove concavity on the open set $\{0<Q<I\}$, where all four probabilities are positive: $p_0=\det(I-Q)>0$ and $p_{12}=\det Q>0$, so $p_1p_2>0$ by \eqref{eq:cov}, and $p_1+p_2=\lambda_1+\lambda_2-2\lambda_1\lambda_2>0$ for the eigenvalues $\lambda_1,\lambda_2\in(0,1)$ of $Q$. On that set concavity is equivalent to $f''(0)\le0$ for every $Q$ and every Hermitian $D$, where $f(t)=\Ent(Q+tD)$.

Put $\varphi(x)=-x\log x$ and
\[
G(a,d,g)=\varphi(g)+\varphi(a-g)+\varphi(d-g)+\varphi(1-a-d+g),
\]
so that $\Ent(Q)=G(a,d,\det Q)$ by \eqref{eq:p}. The function $G$ is a sum of concave functions of affine maps, hence concave, and
\begin{equation}\label{eq:L}
\frac{\partial G}{\partial g}=\log\frac{p_1p_2}{p_0p_{12}}=:L\;\ge0
\end{equation}
by \eqref{eq:cov}. Along $Q+tD$ the entries $a,d$ move affinely, $\frac{d}{dt}\det(Q+tD)\big|_{0}=\tr(\adj(Q)D)=:\gamma$ and $\frac{d^2}{dt^2}\det(Q+tD)=2\det D$. Writing $x=D_{11}$, $y=D_{22}$, the chain rule gives
\begin{equation}\label{eq:f2}
f''(0)=-S+2L\det D,\qquad S=\frac{\ell_{12}^2}{p_{12}}+\frac{\ell_1^2}{p_1}+\frac{\ell_2^2}{p_2}+\frac{\ell_0^2}{p_0},
\end{equation}
with $\ell_{12}=\gamma$, $\ell_1=x-\gamma$, $\ell_2=y-\gamma$, $\ell_0=\gamma-x-y$. One checks directly that
\begin{equation}\label{eq:id}
\ell_1\ell_2-\ell_{12}\ell_0=xy .
\end{equation}

\begin{lemma}\label{lem:adj}
For $0<Q<I$ on $\mathbb C^2$, the matrix $\adj Q$ has the same eigenvalues as $Q$; in particular $0\le\adj Q\le I$.
\end{lemma}

\begin{proof}
$\adj Q=(\det Q)\,Q^{-1}$ has eigenvalues $\lambda_1\lambda_2/\lambda_1=\lambda_2$ and $\lambda_1\lambda_2/\lambda_2=\lambda_1$.
\end{proof}

\begin{lemma}\label{lem:A}
Let $0<Q<I$ and $D\ge0$ on $\mathbb C^2$, and put $A=\gamma\,(\tr D-\gamma)=-\ell_{12}\ell_0$ and $r=\sqrt{p_0p_{12}}$. Then $A\ge(p_1+p_2+2r)\det D$.
\end{lemma}

\begin{proof}
Work in an orthonormal eigenbasis of $\adj Q$, with eigenvalues $\mu_1=\lambda_2$, $\mu_2=\lambda_1$ (Lemma~\ref{lem:adj}), and let $\alpha,\delta\ge0$ be the diagonal entries of $D$ in that basis, so $\alpha\delta\ge\det D$. Then $\gamma=\mu_1\alpha+\mu_2\delta$ and $\tr D-\gamma=(1-\mu_1)\alpha+(1-\mu_2)\delta$, and
\[
A=\mu_1(1-\mu_1)\alpha^2+\mu_2(1-\mu_2)\delta^2+\bigl[\mu_1(1-\mu_2)+\mu_2(1-\mu_1)\bigr]\alpha\delta .
\]
The bracket equals $\lambda_1+\lambda_2-2\lambda_1\lambda_2=p_1+p_2$, and by the arithmetic--geometric mean inequality the first two terms are at least $2\sqrt{\mu_1\mu_2(1-\mu_1)(1-\mu_2)}\,\alpha\delta=2r\,\alpha\delta$, because $\lambda_1\lambda_2=p_{12}$ and $(1-\lambda_1)(1-\lambda_2)=p_0$. Hence $A\ge(p_1+p_2+2r)\alpha\delta\ge(p_1+p_2+2r)\det D$.
\end{proof}

\begin{proof}[Proof of Theorem~\ref{thm:main}]
Fix $0<Q<I$ and a Hermitian $D$. If $\det D\le0$ then $f''(0)\le0$ by \eqref{eq:f2} and \eqref{eq:L}. Otherwise $D$ is definite, and since \eqref{eq:f2} is even in $D$ we may assume $D\ge0$. Put $s=\sqrt{p_1p_2}$ and $r=\sqrt{p_0p_{12}}$, so $s\ge r$ by \eqref{eq:cov}. Applying $u^2/p+v^2/q\ge2|uv|/\sqrt{pq}$ to the pairs $(\ell_{12},\ell_0)$ and $(\ell_1,\ell_2)$,
\[
S\ \ge\ \frac{2|\ell_{12}\ell_0|}{r}+\frac{2|\ell_1\ell_2|}{s}=\frac{2\bigl(|A|s+|B|r\bigr)}{rs},\qquad B:=\ell_1\ell_2 .
\]
The logarithmic mean of $s^2$ and $r^2$ is at least their geometric mean, so $L=\log(s^2/r^2)\le(s^2-r^2)/(rs)$. By \eqref{eq:f2} it therefore suffices to show
\begin{equation}\label{eq:star}
|A|s+|B|r\ \ge\ (s^2-r^2)\det D .
\end{equation}
By Lemma~\ref{lem:adj} and $D\ge0$ we have $\gamma\ge0$ and $\tr D-\gamma=\tr((I-\adj Q)D)\ge0$, so $A\ge0$; and $A+B=xy\ge\det D\ge0$ by \eqref{eq:id}.

If $B<0$, then $A>xy\ge\det D$, and $s\le\frac12(p_1+p_2)\le1$ gives $As\ge s^2\det D\ge(s^2-r^2)\det D$.

If $B\ge0$, then $As+Br=(s-r)A+r\,xy$. Lemma~\ref{lem:A} and $p_1+p_2\ge2s$ give $A\ge(2s+2r)\det D\ge(s+r)\det D$, hence $As+Br\ge(s-r)(s+r)\det D+r\,xy\ge(s^2-r^2)\det D$.

In both cases \eqref{eq:star} holds, so $f''(0)\le0$. Thus $\Ent$ is concave on $\{0<Q<I\}$ and, by continuity, on $\{0\le Q\le I\}$.
\end{proof}

\begin{remark}\label{rem:sharp}
The bound is far from tight in the middle of the parameter range. Over $2\times10^5$ random pairs $(Q,D)$ with $D$ definite, the ratio $2L\det D/S$ never exceeded $0.28$. Equality $f''(0)=0$ occurs, for instance, at $b=0$, $a=d$ (that is, $Q=cI$) with $D$ off-diagonal: there $P_Q$ is a product measure, and moving $Q$ off the diagonal keeps the one-point marginals fixed while $P_Q$ changes only at second order, so the entropy, which is maximal among laws with those marginals, changes at fourth order.
\end{remark}

\section{The second variation for general $E$}\label{sec:general}

For $T\subseteq E$ write $m_T=\det Q_T$ (with $m_\emptyset=1$). Inclusion--exclusion gives
\begin{equation}\label{eq:mob}
P_Q[Y=S]=\sum_{T\supseteq S}(-1)^{|T\setminus S|}m_T ,
\end{equation}
so $P_Q[Y=\cdot\,]$ is an affine function of the vector of principal minors. Differentiating $\Ent=\sum_S\varphi(P_Q[Y=S])$ twice along $Q+tD$ gives the following.

\begin{proposition}\label{prop:second}
Let $0<Q<I$, $D$ Hermitian, $p_S=P_Q[Y=S]$, and $m_T'$, $m_T''$ the first two derivatives of $\det(Q+tD)_T$ at $t=0$. Then
\[
\frac{d^2}{dt^2}\Ent(Q+tD)\Big|_{t=0}=-\sum_{S\subseteq E}\frac{1}{p_S}\Bigl(\sum_{T\supseteq S}(-1)^{|T\setminus S|}m_T'\Bigr)^{2}+\sum_{|T|\ge2}\Lambda_T\,m_T'',
\]
where $\Lambda_T=-\sum_{S\subseteq T}(-1)^{|T\setminus S|}\log p_S$.
\end{proposition}

\begin{proof}
By \eqref{eq:mob}, $\frac{d}{dt}p_S=\sum_{T\supseteq S}(-1)^{|T\setminus S|}m_T'$ and $\frac{d^2}{dt^2}p_S=\sum_{T\supseteq S}(-1)^{|T\setminus S|}m_T''$. Since $\varphi'(x)=-1-\log x$ and $\varphi''(x)=-1/x$,
\[
\frac{d^2}{dt^2}\Ent=\sum_S\varphi''(p_S)\Bigl(\frac{dp_S}{dt}\Bigr)^2+\sum_T m_T''\sum_{S\subseteq T}(-1)^{|T\setminus S|}\varphi'(p_S).
\]
The constant $-1$ in $\varphi'$ contributes $\sum_{S\subseteq T}(-1)^{|T\setminus S|}=0$ for $T\ne\emptyset$, and $m_T''=0$ for $|T|\le1$ because $m_T$ is then affine in $t$.
\end{proof}

When $Q<I$, the measure $P_Q$ is the $L$-ensemble with $L=Q(I-Q)^{-1}$, that is $p_S=\det L_S/\det(I+L)$. The normalisation cancels in $\Lambda_T$ for $T\ne\emptyset$, so
\[
\Lambda_T=-\sum_{S\subseteq T}(-1)^{|T\setminus S|}\log\det L_S .
\]
For $T=\{i,j\}$,
\[
\Lambda_{ij}=\log\frac{L_{ii}L_{jj}}{\det L_{\{i,j\}}}\ \ge0
\]
by the Hadamard--Fischer inequality; for $|E|=2$ this is \eqref{eq:L}. The proof of Theorem~\ref{thm:main} uses exactly two features of the case $|E|=2$: the only curvature term is $\Lambda_{12}\,m_{12}''=2\Lambda_{12}\det D$, with $\Lambda_{12}\ge0$, and $\Lambda_{12}$ is controlled by the quadratic form through the logarithmic-mean inequality. For $|E|\ge3$ the coefficients $\Lambda_T$ with $|T|\ge3$ take both signs: for $|E|=3$, in $3000$ random kernels with spectra in $[0.05,0.95]$, $\Lambda_{123}$ was negative in $1508$ cases, and $\Lambda_T m_T''>0$ occurred for both $|T|=2$ and $|T|=3$. The curvature terms therefore have to be compared with the quadratic form jointly.

\section{Numerical evidence for three and four points}\label{sec:numerics}

We searched for counterexamples in two ways, with kernels $Q=U\operatorname{diag}(\lambda)U^*$ for $U$ from the QR factorisation of a Gaussian matrix (real or complex) and $\lambda\in(0,1)^n$.

\emph{Midpoint gaps.} For $n=|E|=3$ (real) and $n=4$ (complex) we minimised $\Ent\bigl((Q_1+Q_2)/2\bigr)-\bigl(\Ent(Q_1)+\Ent(Q_2)\bigr)/2$ by Nelder--Mead from the worst of $2000$ and $3000$ random pairs. The minima were $8\times10^{-15}$ and $2\times10^{-8}$, attained as $Q_1-Q_2\to0$.

\emph{Second derivatives.} We maximised the central difference $\bigl(\Ent(Q+hD)+\Ent(Q-hD)-2\Ent(Q)\bigr)/h^2$ with $h=10^{-3}$ and $\|D\|_F=1$, keeping the spectrum of $Q$ at distance at least $5h$ from $\{0,1\}$. For $n=3$ (real) and $n=4$ (complex) the maxima were $-2\times10^{-6}$ and $-1.2\times10^{-6}$, approached only near scalar kernels $Q=cI$ (in every optimised kernel the eigenvalues agreed to within $4\times10^{-4}$). Near those points, for $n=3$, $Q=cI+\varepsilon A$ with $c\in[0.2,0.8]$, $\varepsilon$ down to $3\times10^{-3}$ and $400$ random real symmetric pairs $(A,D)$, half of them with $D$ off-diagonal, the second derivative stayed negative. This matches Remark~\ref{rem:sharp}: scalar kernels give product measures, and off-diagonal perturbations preserve the one-point marginals.

The formula of Proposition~\ref{prop:second} agrees with central differences to relative accuracy $3\times10^{-6}$ at $n=3$.

\section*{Acknowledgements and funding}
The research presented in this paper was supported by the European Research
Council (ERC) under the European Union's Horizon Europe research and
innovation programme (grant agreement No.~101041711), by the Simons
Foundation as part of the Collaboration on the Mathematical and Scientific
Foundations of Deep Learning, by Heights Labs, by Convex Nexus Capital, by
the Israel Science Foundation (grant number 2258/19), and by the Israel
Science Foundation (ISF Grant 4101/25). Computational resources were
provided in part by Google Cloud through its education credits program
(credit 530550103).

\section*{Declaration of competing interest}
The author holds equity in Heights Labs and in Convex Nexus Capital and
has a management role at Convex Nexus Capital. These entities are included
in the funding declaration above.

\section*{Declaration of generative AI and AI-assisted technologies}
The paper comes out of Math Brain, an ongoing project on the use of AI systems for research mathematics. In November 2026 Claude Opus~5.5 (Anthropic), working in Claude Code, found the proof of Theorem~\ref{thm:main} and the formula of Proposition~\ref{prop:second}, wrote and ran the computations of Section~\ref{sec:numerics}, and wrote the paper. The author directed the work and is responsible for the content of the paper and for decisions concerning its dissemination.

\section*{Data availability}
No data were used. The code for the computations of Remark~\ref{rem:sharp} and Sections~\ref{sec:general}--\ref{sec:numerics}, with their outputs, is available at\\ \url{https://github.com/yspennstate/lyons-determinantal-entropy-two-points}.

\begin{thebibliography}{9}

\bibitem{Ly03}
R.~Lyons, Determinantal probability measures, \emph{Publ. Math. Inst. Hautes \'Etudes Sci.} \textbf{98} (2003), 167--212.

\bibitem{Ly14}
R.~Lyons, Determinantal probability: basic properties and conjectures, in \emph{Proceedings of the International Congress of Mathematicians, Seoul 2014}; arXiv:1406.2707.

\bibitem{LS03}
R.~Lyons and J.~E.~Steif, Stationary determinantal processes: phase multiplicity, Bernoullicity, entropy, and domination, \emph{Duke Math. J.} \textbf{120} (2003), 515--575.

\end{thebibliography}
\end{document}
